% !TeX program = xelatex
% !TeX encoding = UTF-8
% =====================================================================
%  通用中文技术文档 / 学习笔记模板
%  编译方式：xelatex（中文必须用 XeLaTeX，不要用 pdfLaTeX）
%  推荐命令：latexmk -xelatex main.tex
%  手动编译：xelatex main -> biber main -> xelatex main -> xelatex main
% =====================================================================
\documentclass[11pt,a4paper]{article}

% ---------------------------------------------------------------------
% 1. 中文支持（ctex 会自动检测 Windows / macOS / Linux 的中文字体）
%    如编译报字体错误，可手动指定：
%      Windows : \usepackage[fontset=windows]{ctex}
%      macOS   : \usepackage[fontset=mac]{ctex}
%      Linux   : \usepackage[fontset=fandol]{ctex}
% ---------------------------------------------------------------------
\usepackage{ctex}

% ---------------------------------------------------------------------
% 2. 版面设置
% ---------------------------------------------------------------------
\usepackage[margin=2.5cm]{geometry}   % 页边距
\usepackage{setspace}
\onehalfspacing                        % 1.3~1.5 倍行距，中文更易读
\usepackage{parskip}                  % 段间留空、段首不缩进（技术文档风格）
\usepackage{xcolor}

% ---------------------------------------------------------------------
% 3. 常用宏包
% ---------------------------------------------------------------------
\usepackage{graphicx}                 % 插图
\usepackage{booktabs}                 % 三线表（\toprule \midrule \bottomrule）
\usepackage{array}                    % 增强表格列格式
\usepackage{enumitem}                 % 列表间距微调
\usepackage{amsmath,amssymb,amsthm}   % 数学公式与符号
\usepackage[hidelinks]{hyperref}      % 交叉引用（hidelinks：先加载，后面统一着色）

% ---------------------------------------------------------------------
% 4. 定理类环境（按需使用，不用也不影响）
%    编号带章节号，如“定理 2.1”
% ---------------------------------------------------------------------
\newtheorem{theorem}{定理}[section]
\newtheorem{lemma}[theorem]{引理}
\newtheorem{proposition}[theorem]{命题}
\theoremstyle{definition}
\newtheorem{definition}[theorem]{定义}
\theoremstyle{remark}
\newtheorem*{remark}{注}

% ---------------------------------------------------------------------
% 5. 代码排版（listings，无需安装额外软件）
%    若追求更漂亮的高亮，可改用 minted（需安装 Python + Pygments，见 README）
%    注意：listings 对中文支持较差，代码内的注释/字符串建议用英文；
%          正文里讲代码不受影响。
% ---------------------------------------------------------------------
\definecolor{codebg}{RGB}{248,248,248}
\definecolor{framegray}{RGB}{210,210,210}
\definecolor{kwblue}{RGB}{0,0,180}
\definecolor{cmtgreen}{RGB}{0,128,0}
\definecolor{strred}{RGB}{163,21,21}

\usepackage{listings}
\lstset{
  backgroundcolor=\color{codebg},     % 浅灰背景
  basicstyle=\ttfamily\footnotesize,  % 等宽字体
  keywordstyle=\color{kwblue}\bfseries,
  commentstyle=\color{cmtgreen}\itshape,
  stringstyle=\color{strred},
  numbers=left,                       % 左侧行号
  numberstyle=\tiny\color{gray},
  numbersep=8pt,
  frame=single,                       % 边框
  rulecolor=\color{framegray},
  framesep=4pt,
  breaklines=true,                    % 长行自动换行
  showstringspaces=false,
  tabsize=4,
  captionpos=b,                       % 代码标题在下方
  columns=fixed,
  extendedchars=true,
  upquote=true,
}
\renewcommand{\lstlistingname}{代码}   % 代码块标题中文化：代码 1, 代码 2, ...
% 常用语言快捷切换：\begin{lstlisting}[language=Python] ...
%   内置支持 Python / Java / C / C++ / Bash / SQL / Matlab 等

% ---------------------------------------------------------------------
% 6. 参考文献：GB/T 7714-2015 国标格式（顺序编码制）
%    文献条目写在 references.bib 中
% ---------------------------------------------------------------------
\usepackage[backend=biber,style=gb7714-2015]{biblatex}
\addbibresource{references.bib}

% ---------------------------------------------------------------------
% 7. 超链接配色（放在所有宏包之后）
% ---------------------------------------------------------------------
\hypersetup{
  colorlinks=true,
  linkcolor=structureblue,
  citecolor=structureblue,
  urlcolor=structureblue,
  bookmarksnumbered=true,
}
\definecolor{structureblue}{RGB}{26,82,118}

% ---------------------------------------------------------------------
% 8. 文档信息：每次写新报告只改这里
% ---------------------------------------------------------------------
\title{\textbf{在这里写报告标题}}
\author{张三 \\
  软件工程专业 \,$\cdot$\, 学号 20230000 \\
  \texttt{zhangsan@example.com}}
\date{\today}   % \today 自动取编译当天日期；也可手写成 2026年9月14日

\begin{document}

\maketitle

% ==================== 正文从这里开始 ====================

\section{概述}

这是一个通用的中文技术文档模板，适用于课程实验报告、技术总结和学习笔记。
本文同时演示了模板支持的各种排版元素，写完后把示例内容删掉即可。

行内数学公式用单个美元符号包裹，例如 $E = mc^2$、$f(x)=\int_0^x g(t)\,\mathrm{d}t$。
常用的排版结构包括：

\begin{itemize}[itemsep=2pt,topsep=4pt]
  \item 无序列表：要点一、要点二；
  \item 有序列表见下方；
  \item 以及对外部资料的引用，如文献\cite{knuth1997taocp}。
\end{itemize}

\begin{enumerate}[itemsep=2pt,topsep=4pt]
  \item 第一步，复制整个模板文件夹；
  \item 第二步，修改标题、作者等文档信息；
  \item 第三步，删除示例，开始写自己的内容。
\end{enumerate}

\section{数学公式与定理}

行间公式用 equation 环境并自动编号，方便在正文里交叉引用：
\begin{equation}
  \label{eq:gaussian}
  f(x) = \frac{1}{\sqrt{2\pi}\,\sigma}
         \exp\!\left(-\frac{(x-\mu)^2}{2\sigma^2}\right).
\end{equation}

公式\eqref{eq:gaussian}是高斯分布的概率密度函数。多行对齐用 align 环境：
\begin{align}
  (a+b)^2 & = a^2 + 2ab + b^2, \\
  (a-b)^2 & = a^2 - 2ab + b^2.
\end{align}

定理、定义等环境的用法如下。

\begin{definition}[勾股定理]
  设直角三角形两条直角边长度为 $a$、$b$，斜边长度为 $c$，则
  \begin{equation*}
    a^2 + b^2 = c^2.
  \end{equation*}
\end{definition}

\begin{theorem}
  \label{thm:example}
  任意有限状态的马尔可夫链至少存在一个平稳分布。
\end{theorem}

\begin{proof}
  证明从略，这里只演示环境样式。
\end{proof}

\section{图表}

\subsection{图片}

把图片文件放进 \texttt{figures/} 目录，然后用 \texttt{$\backslash$includegraphics} 插入。
下面先用方框作占位演示（实际使用时替换成注释中的命令）：

\begin{figure}[htbp]
  \centering
  \fbox{\parbox[c][4cm][c]{0.62\linewidth}{\centering 图片占位框\\[6pt]
    {\footnotesize 实际使用时替换为下面的命令：}\\
    {\footnotesize\tt $\backslash$includegraphics\{figures/demo.png\}}}}
  \caption{图片标题放在图的下方}
  \label{fig:demo}
\end{figure}

如图~\ref{fig:demo} 所示，建议优先使用矢量图（PDF/SVG 转 PDF），放大不模糊。

\subsection{三线表}

正式文档推荐三线表，只保留三条横线，不要竖线，如表~\ref{tab:example}。

\begin{table}[htbp]
  \centering
  \caption{表格标题放在表的上方}
  \label{tab:example}
  \begin{tabular}{lcc}
    \toprule
    算法 & 时间复杂度 & 空间复杂度 \\
    \midrule
    冒泡排序 & $O(n^2)$ & $O(1)$ \\
    归并排序 & $O(n\log n)$ & $O(n)$ \\
    快速排序 & $O(n\log n)$（平均） & $O(\log n)$ \\
    \bottomrule
  \end{tabular}
\end{table}

\section{代码}

行内代码用 \texttt{lstinline}，例如 \lstinline{int x = 0;}。
多行代码用 lstlisting 环境，指定语言即可获得关键字高亮和行号：

\begin{lstlisting}[language=Python, caption={Python 示例：二分查找}]
def binary_search(arr, target):
    lo, hi = 0, len(arr) - 1
    while lo <= hi:
        mid = (lo + hi) // 2
        if arr[mid] == target:
            return mid
        elif arr[mid] < target:
            lo = mid + 1
        else:
            hi = mid - 1
    return -1
\end{lstlisting}

切换语言只需改 \texttt{language=Java}、\texttt{language=C++}、\texttt{language=Bash} 等。

\section{参考文献的写法}

参考文献条目统一写在 \texttt{references.bib} 里，正文用
\texttt{$\backslash$cite\{key\}} 引用，文末自动生成国标 GB/T 7714-2015
格式的参考文献表，例如\cite{tanenbaum2014modern}和中文文献\cite{yan2008wcdd}。
引用、编号、排序全部自动处理，增删文献无需手动调整编号。

% ==================== 参考文献 ====================
\printbibliography[heading=bibintoc,title=参考文献]

\end{document}
